% page 332 du fac-simile (image p-331 du scan)
% bloc 152.4 x 229.0 mm ; marge gauche 8.0 mm
\pgc{-12.700mm}{0.000mm}{
\l{\cel{3}324}
\l{}
\l{\sou{kursoro}. (\sou{tekn}.)Bloko plata, mobila alonge linealo,}
\l{\cel{3}limbo gradizita, e c., uzata kom indexo. EFIS.}
\l{}
\l{\sou{kurta}. I.Di qua la extenseso, de un de sua extremaji til}
\l{\cel{3}la altra, sinse di longeso o di alteso, esas poka.}
\l{\cel{3}II. Di qua la duro esas poka. - DEFIS.}
\l{}
\l{\sou{kurtajar}. (\sou{netrans.}) Agar kom mediacero, en komerco, in-}
\l{\cel{3}ter la vendanto e la kompranto, po pago.- FRS.}
\l{}
\l{\sou{kurteno}. I. Peco de stofo tensita por celar o shirmar ulo}
\l{\cel{3}od ulu. - II. (\sou{Teatro}). Kurteno qua celas la +stejo}
\l{\cel{3}kande ne pleesas. - III. Facio di la muro di remparo}
\l{\cel{3}inter du bastioni. - EFIS.}
\l{}
\l{\sou{kurtezar}. (\sou{trans}.)I. Prizentar (ad ulu) sua homaji por}
\l{\cel{3}obtenar lua favoro ; esforcar obtenar favoreso (da o}
\l{\cel{3}da ulu, precipue da muliero) per anuncar o dicar a lu}
\l{\cel{3}ulo qua interesas lu ed efikas a lua vanitato. - DEFIS.}
\l{}
\l{\sou{kurva}. Di qua omna elemento su eskartas, per deviaco pres-}
\l{\cel{3}ke neperceptebla, e segun +leyo (lego) fixa, de la rek-}
\l{\cel{3}to o de la plano. - (\sou{geom}.) Lineo qua esas nek rekta,}
\l{\cel{3}nek kompozuro ek rekti.\cel{2}- DEFIS.}
\l{}
\l{\sou{kurvimetro}.\cel{2}Instrumento por mezurar la longeso di la}
\l{\cel{3}kurvi, sur papero. - DEFIS.}
\l{}
\l{\sou{kuseno}. Mikra sako quadrata, ek stofo, ledro, polsteriz-}
\l{\cel{3}ita ye plumi, krini, e c., por ke onu povez apogar su}
\l{\cel{3}an olu. - EFIS.}
\l{}
\l{\sou{kushar}. (\sou{trans., sur}). I. Pozar aden lito, instalar en lito.}
\l{\cel{3}- II. Pozar aden lito, instalar en lito, por repozar}
\l{\cel{3}nokto-dorme. - F.}
\l{}
\l{\sou{kuskuto}. (\sou{bot.)}\cel{3}Herbo parazita filoforma, qua plekto-}
\l{\cel{3}kaptas la planti vicina e sufokigas li. - L.\cel{1}\sou{cuscuta.}}
\l{\cel{3}- FISL.}
\l{}
\l{\sou{kuspido}. (geom.)\cel{3}Punto ube kurvo su dividas aden du ben-}
\l{\cel{3}di,kun tangento komuna a ta punto. -}
\l{}
\l{\sou{kustar.}\cel{1}(trans.)I. (\sou{aludante objekto, laboro)}. Igar necesa}
\l{\cel{3}la pago di pekunio-quanto por obtenesar kambie. - II.}
\l{\cel{3}(\sou{metaf.)}\cel{1}Igar necesa sakrifiko por obtenesor. - III.}
\l{\cel{3}Havar kom konsequajo ulo quo impozas spenso. - DEFIS.}
\l{}
\l{\sou{kustumar.}\cel{1}(trans.)(aludante populo, naciono od individuo).}
\l{\cel{3}Agar konforme a la maniero establisita da longa experien-}
\l{\cel{3}co o praktiko. - EFIS.}
}
